% Compile with: latexmk -lualatex -pdf corpus-stats.tex
\documentclass[11pt,a4paper]{article}

% ---------- Fonts & Languages ----------
\usepackage{fontspec}
% \setmainfont{Charis SIL}

\usepackage{polyglossia}
\setmainlanguage{english}

% ---------- Typography ----------
\usepackage{microtype}

% ---------- Math, statistics & tables ----------
\usepackage{amsmath,amssymb}
\usepackage{siunitx}
\usepackage{booktabs}

% ---------- Plots ----------
\usepackage{pgfplots}
\pgfplotsset{compat=1.18}

% ---------- Layout ----------
\usepackage{geometry}
\geometry{margin=2.5cm}
\usepackage{hyperref}
\hypersetup{colorlinks=true, linkcolor=blue, urlcolor=blue}

\title{Corpus Linguistics in LuaLaTeX \\ \large Statistics, Tables, and Zipf Plots}
\author{Your Name}
\date{\today}

\begin{document}
\maketitle

This worked example shows how to produce publication-ready statistical
output directly in LuaLaTeX: aligned decimal tables with \texttt{siunitx},
and vector-based plots with \texttt{pgfplots}. No Python screenshots required
--- everything stays in the \texttt{.tex} source and recompiles when the data
change.

\section{Aligned Statistical Tables with \texttt{siunitx}}

The \texttt{S} column type in \texttt{siunitx} aligns numbers on their decimal
point. Scientific notation is formatted automatically, and the \textit{p} in
\textit{p}-value stays correctly italicized.

\begin{center}
\begin{tabular}{l S[table-format=2.3] S[table-format=1.3e-1]}
\toprule
\textbf{Language} & \textbf{Mean Word Length} & \textbf{\textit{p}-value} \\
\midrule
French  & 4.812 & 3.4e-5 \\
Spanish & 5.123 & 1.2e-6 \\
Italian & 5.401 & 4.5e-4 \\
\bottomrule
\end{tabular}
\end{center}

In running text you can also use \verb|\num{1.23e-4}| to format a number as
$\num{1.23e-4}$, and \verb|\SI{95}{\percent}| for a value with its unit:
$\SI{95}{\percent}$.

\section{Zipf's Law Plot with \texttt{pgfplots}}

The plot below is generated purely in LaTeX from inline coordinates. The
output is vector-based, so it stays sharp at any zoom level and the fonts
match the rest of the document.

\begin{center}
\begin{tikzpicture}
\begin{loglogaxis}[
    xlabel={Rank (\(r\))},
    ylabel={Frequency (\(f\))},
    title={Zipf's Law in a Romance Corpus},
    grid=major,
    width=10cm, height=7cm,
    legend pos=south west,
]
\addplot[only marks, mark=*, blue] coordinates {
    (1, 1000) (2, 500) (3, 333) (4, 250) (5, 200)
    (10, 100) (20, 50) (50, 20) (100, 10)
};
\addplot[red, thick, domain=1:100] {1000/x};
\legend{Observed, $f = 1000/r$}
\end{loglogaxis}
\end{tikzpicture}
\end{center}

The red curve is the theoretical Zipf law $f = C/r$ for comparison. Replace
the inline coordinates with your own corpus data when you have them.

\section{Loading Data from a CSV}

For a real corpus you do not want to retype counts into the \TeX{} source.
Load them from a CSV with \texttt{pgfplotstable}:

\begin{verbatim}
\usepackage{pgfplotstable}
% ...
\addplot table [x=rank, y=freq, col sep=comma] {my_corpus.csv};
\end{verbatim}

The CSV needs a header row (e.g. \texttt{rank,freq}) and one data row per
observation. Every recompile re-reads the file, so updating the corpus
updates the plot without touching the LaTeX source.

\section{Short Reference}

\begin{itemize}
    \item \verb|\num{1.23e-4}| --- format a number in text.
    \item \verb|S[table-format=2.3]| --- column type that aligns decimals;
          \texttt{2.3} means up to 2 digits before and 3 after.
    \item \verb|\SI{95}{\percent}| --- number with a unit.
    \item \verb|\begin{loglogaxis}| --- use instead of \texttt{axis} for log-log plots.
    \item \verb|\addplot table [x=..., y=...]| --- read columns from a CSV.
\end{itemize}

\end{document}
